\documentclass[12pt]{article}
\usepackage{amsmath}
\usepackage{mathrsfs}
\usepackage{eucal}
\usepackage{amssymb}
\usepackage[french]{babel}
\usepackage[latin1]{inputenc}
\usepackage[T1]{fontenc}
\usepackage[all]{xy}
%\usepackage{graphicx}
%\usepackage{epsf}

\begin{document}

\hfuzz 8 pt
\parindent = 0 cm

\centerline{ \textsc{ \large Universit� de Sherbrooke}}
\centerline{\textsc{ \normalsize D�partement de math�matiques}}


\bigskip

\centerline{\large\bf MAT417 -- M�thodes num�riques en alg�bre
lin�aire}

\bigskip
\centerline{\Large\bf Laboratoire 4}
\bigskip


\begin{description}

\item[1)] Consid�rons le syst�me d'�quations lin�aires

$\left( \begin{array}{rrr} -2 & 1 & 1 \\ 1 & 3 & 2 \\ 1 & 2 & 1
\end{array} \right) \left(\begin{array}{r} x_1 \\ x_2 \\ x_3
\end{array} \right) = \left(\begin{array}{r} 2 \\ 9 \\ 6
\end{array} \right)$

    \begin{description}

    \item[a)] Effectuer une d�composition de Choleski de la matrice des coefficients du syst�me.

    \item[b)] R�soudre le syst�me obtenu.

    \end{description}


\item[2)] Soit $x \in \mathbb{R}^n$ pour un $n>0$ et $\| x \|_p =
\left( \sum\limits_{i=1}^{n} |x_i|^p \right)^{\frac{1}{p}}$ la
norme $p$ de $x$.

    \begin{description}

    \item[a)] Montrer que $\| x \|_{\infty} \leqslant \| x \|_p$
    pour toute valeur de $p$.

    \item[b)] Soit $A$ une matrice et $\| A \|_{\infty}$ la norme
    matricielle induite de la norme vectorielle. Montrer que
    $\| A \|_{\infty} = \max\limits_{i=1,2,\ldots,n} \left( \sum\limits_{j=1}^{n} | a_{ij} | \right)$.

    \item[c)]  Soit $A$ une matrice et $\| A \|_{1}$ la norme
    matricielle induite de la norme vectorielle. Montrer que
    $\| A \|_{1} = \max\limits_{j=1,2,\ldots,n} \left( \sum\limits_{i=1}^{n} | a_{ij} | \right)$.

    \end{description}


\item[3)] Soit $A$ une matrice $n$ par $n$ pour un certain $n>0$
et $x \in \mathbb{R}^{n}$. Consid�rons la norme de Frobenius
d�finie par

$$\| A \|_{F} = \sqrt{\sum\limits_{i=1}^{n}\sum\limits_{j=1}^{n} a_{ij}^2}.$$

    \begin{description}

    \item[a)] Montrer que $\| A \|_{F}$ est bel et bien une
    norme, c'est-�-dire que pour toutes matrices $A$ et $B$ et
    pour toute constante $\alpha$

        \begin{description}

        \item[i)] $\| A \|_{F} \geqslant 0$

        \item[ii)] $\| A \|_{F} = \iff A=0 $

        \item[iii)] $\| \alpha A \|_{F} = | \alpha | \cdot \| A
        \|_{F}$

        \item[iv)] $\| A + B \|_{F} \leqslant \| A \|_{F} + \| B \|_{F}$

        \end{description}

    \item[b)] Montrer que la norme de Frobenius est compatible
    avec la norme vectorielle $\| \cdot \|_2$, c'-est-�-dire que
    $\| Ax \|_2 \leqslant \| A \|_F \cdot \| x \|_2$.

    \item[c)] Montrer que $\| A \|_2 \leqslant \| A \|_F$.

    \item[d)] D�terminer une matrice $A$ telle que $\| A \|_2 \neq \| A \|_F$.

    \item[e)] En d�duire que $\| A \|_{F}$ ne peut �tre une norme
    induite d'une norme vectorielle, c'est-�-dire que $\| A \|_F
    \neq \sup\limits_{y \in \mathbb{R}^{n}_{*}} \frac{\| Ay \|}{\| y
    \|}$

    \end{description}


\item[4)] Montrer que ${\rm cond}(A) \geqslant 1$ pour toute
matrice $A$.



\item[5)] Consid�rons la matrice $\left( \begin{array}{rrr} 3 & 1
& 6 \\ 2 & 1 & 3 \\ 1 & 1 & 1 \end{array} \right).$

    \begin{description}

    \item[a)] Calculer ${\rm cond}(A)$.

    \item[b)] Conditionner les lignes de la matrice $A$ avec la
    norme $\| \cdot \|_{1}$, puis calculer le conditionnement de la nouvelle
    matrice.

    \item[c)] Conditionner les colonnes de la matrice r�sultante avec la
    norme $\| \cdot \|_{1}$, puis calculer le conditionnement de la nouvelle
    matrice.

    \item[d)] Refaire les parties b) et c), mais en normalisant
    cette fois avec la norme $\| \cdot \|_{\infty}$.

    \item[e)] En r�solvant le syst�me d'�quation

$$\left( \begin{array}{rrr} 3 & 1 & 6 \\ 2 & 1 & 3 \\ 1 & 1 & 1
\end{array} \right)
\left(\begin{array}{r} x_1 \\ x_2 \\ x_3
\end{array} \right) = \left(\begin{array}{r} 2 \\ 7 \\ 4
\end{array} \right),$$

on a obtenu en arithm�tique � $m=3$ chiffres significatifs la
solution approximative

$$ \left(\begin{array}{r} x_1 \\ x_2 \\ x_3
\end{array} \right) = \left(\begin{array}{r} 18,8 \\ -6,78 \\
-8,01
\end{array} \right).$$

D�terminer des bornes pour l'erreur d'approximation relative de
cette solution.

    \end{description}


\item[6)] Consid�rons le syst�me it�ratif

$$x^{(k+1)} = F x^{(k)} + c$$ o�

$$F = \left( \begin{array}{rrr} 0 & 0 & 2 \\ 1 & 1/2 & -4 \\
6 & 3 & 1/4
\end{array} \right).$$

Montrer que ce syst�me converge pour toute valeur de $x^{(0)}$.

\end{description}

\end{document}
